\section{An unbounded dyadic source proof pays mixed proper heads from first missing twenty onward}
This theorem uses fourteen scoped standard-axiom Lean reports for the
whole fold, power, indexed-gap and finite-source-sum arguments below.
The actual word, CI, Cauchy and proper-head assembly is the written proof,
not an end-to-end Lean word implementation. General C, general A2 and
the ORIGINAL conjecture remain OPEN.

Set $\rho=3/5$, $b=35/32$, $c=3/2$, $q=b^2/(2\rho)=6125/6144$,
and $v=b/(2\rho)=175/192$. For $t\ge2$ put
\[
 s_t=S_t(\rho)=\tfrac32(1-\rho^t).
\]
All capacities, counts, weights and initial runs below are UNBOUNDED.

\subsection{A uniform positive evaluation bound, including the whole append fold}
For EVERY formal word $1^\beta V$, allowing $\beta=0$, let
$M=\beta+\sum_{t\in V}(t-1)$. Then
\[
                       F_{1^\beta V}(\rho)\le10b^M-125/22<10b^M. \tag{1}
\]
The bound does not require ordered capacities. To prove it, first note
\[
 0\le s_t\le b^{t-1}\quad(t\ge2),\qquad
 s_t\ge147/125=1+22/125\quad(t\ge3).                  \tag{2}
\]
For $t=2,3,4,5$ the upper bounds are four explicit rational inequalities.
For $t\ge6$, use $s_t\le3/2\le b^5\le b^{t-1}$.
The seed $[\beta+1](\rho)\le5/2$ and $b^\beta\ge1$ establish
(1) at the empty capacity word, including $\beta=0$.
If (1) holds before appending $t$, then $F'=1+s_tF$ and $M'=M+t-1$.
For $t\ge3$, the negative offset pays the whole new unit:
$(125/22)(s_t-1)\ge1$. For $t=2$, use $b^M\ge1$ and the exact
positive margin
\[
           10(b-24/25)-1-(125/22)/25=97/880>0.
\]
This proves the ENTIRE arbitrary-length fold, not a sampled range.

\subsection{A count-free first-missing penalty and a uniform power block}
Let $K=\kappa_Y\ge20$ and $n=K-2\ge18$. The LAST $n$ non-one
capacities of the ordered actual arm $Y$ are respectively at least
$2,3,\ldots,n+1$. Their skeleton degree is at least
$D_n=n(n+1)/2$. Write $M_0$ for the earlier prefix's skeleton degree.
Since $s_t\le c$, the whole tail fold satisfies
\[
 F_Y(\rho)+2\le c^n(F_{Y_0}(\rho)+2)
                  \le10c^n b^{M_0}.                  \tag{3}
\]
The last bound RETAINS the negative offset $125/22>2$ in (1).
The exact integer power checks are
\[
 q^{224}\le\tfrac12,\qquad
           2240c^{18}\le b^{171},\qquad c\le b^5.    \tag{4}
\]
They imply, for EVERY integer $m\ge0$ and $n\ge18$,
\[
                 m q^m\le224,\qquad2240c^n\le b^{D_n}. \tag{5}
\]
Indeed write $m=224a+r$, $0\le r<224$. Then
$m q^m\le224(a+1)2^{-a}\le224$ by $a+1\le2^a$.
For the second assertion start at $n=18$, use
$D_{n+1}=D_n+n+1$, and multiply by $b^{n+1}\ge b^5\ge c$.
Thus (5) is a whole unbounded block/triangle induction.

\subsection{The ACTUAL forward slope is dyadically thin outside its strict CI head}
Let $M=\beta+\sum(t-1)$ be the main-box degree of $Y$, and
$B=\max(\beta,\max(t-1))$. Put
\[
                            s_Y=\min(\lfloor M/2\rfloor,M-B).
\]
Then for EVERY $m>s_Y$,
\[
                          m\,\Delta F_Y(m)\le2^m.    \tag{6}
\]
First suppose $M<2m$. Positivity of coefficients, (3), and
$M\ge M_0+D_n$ give
\[
 b^{D_n}\frac{F_Y(m)}{2^m}
        \le10c^n q^m.
\]
Multiplication by $m$ and (5) prove $mF_Y(m)\le2^m$,
which is stronger than (6).

It remains to treat $m>M-B=:O$. The first-missing staircase gives
$O\ge1+(K-3)(K-2)/2\ge154$.
If the largest interval is the last capacity interval $[T]$, delete it
to obtain the actual prefix $Y_0$ with skeleton degree $O$. The exact
forward recurrence gives
\[
 \Delta F_Y(m)=F_{Y_0}(m-1)-F_{Y_0}(m-T-1)\le F_{Y_0}(m-1).
\]
Write $j=m-1\ge O$. From (1) and positivity,
\[
                   \frac{\Delta F_Y(m)}{2^m}\le5v^j.
\]
If instead the initial interval $[\beta+1]$ is largest, write
$F_Y=A+P[\beta+1]$ by positive suffix boxes, where
$A$ is their proper-suffix sum and $P=\prod S_t$.
Then coefficientwise
\[
 (1-z)F_Y=(1-z)A+P(1-z^{\beta+1})\preceq A+P=:F_{0,V}.
\]
This is the allowed zero-initial-run formal fold of skeleton degree $O$,
so (1) gives $\Delta F_Y(m)/2^m\le10v^m$.
The exact base $10\cdot101v^{100}<1$, followed by the decreasing-ratio
induction $v(j+2)\le j+1$ for $j\ge100$, proves
$10(j+1)v^j\le1$ for EVERY $j\ge100$. Since $O\ge154$,
both largest-factor cases prove (6). These cases cover the actual minimum
in $s_Y$, including ties, plateaus and zero/negative slopes.

\subsection{An indexed curvature gap for ANY actual mixed owner}
Let $Z=1^\alpha U$ have $h\ge2$ non-one parts, skeleton degree $D$,
weight $a=\alpha+D+h$, and ordinary reversal $Q=q_Z$.
At a mixed/source-prefix recovery rank let $r=a-j\ge0$,
$Q_r\ge Q_{r-1}$ and
\[
                           C=Q_{r+1}-2Q_r+Q_{r-1}>0.
\]
The frozen primitive-center theorem gives $D\ge2r+2$.
The frozen ordinary Hilbert realization gives
$(r+1)Q_{r+1}\le(h+r+1)Q_r$, including every padding and capacity.
We claim the stronger indexed gap
\[
                    (j+1)(2Q_r-Q_{r-1})\ge4Q_{r+1}.  \tag{7}
\]
For $r\ge2$ it suffices to show
$(j+1)(r+1)\ge4(h+r+1)$, because $2Q_r-Q_{r-1}\ge Q_r$.
We have $j+r\ge2h+1$ and $j\ge h+r+3$.
For $r\ge3$ use
$(h+r+4)(r+1)-4(h+r+1)=h(r-3)+r(r+1)\ge0$.
For $r=2$, split at $h=6$ and use the respective two bounds on $j$.
For $r=1$, use $Q_0=1,Q_1=h+1$,
$2Q_2\le(h+1)(h+2)$ and $j\ge2h$; the remaining positive polynomial
is $2h^2-2h-3\ge1$ at $h\ge2$.
For $r=0$, $Q_{-1}=0,Q_0=1,Q_1=h+1$ and $j\ge2h+1$ give (7)
directly. No first-missing bound or C assumption is used for the owner.

\subsection{A whole dyadic source sum pays the indexed gap}
Write $A=F_Z$, $B=F_Y$, $N=j+1$, and
$E_B=1-z+(2z-1)B$. Its entire prefix through $N$ is nondecreasing
in the mixed region; put $w_s=\Delta E_B(s)\ge0$.
Its first possible birth is at least three. The actual doubling law
$A_{i+1}\le2A_i$ gives, for every used $s$,
\[
                       A_{N-s}\ge2^{2-s}A_{N-2}.
\]
The COMPLETE telescoping identity, including the rank-two correction, is
\[
 4\sum_{s=2}^N2^{-s}w_s
                   =1-4\,2^{-N}\Delta B_N.           \tag{8}
\]
It follows from $B_0=B_1=1$ and $E_B(s)=2B_{s-1}-B_s$ for $s\ge2$.
Consequently
\[
 \Delta(E_BA)_N\ge A_{N-2}(1-4\,2^{-N}\Delta B_N).    \tag{9}
\]
The exact mixed tail increment is $T=C-\Delta(E_BA)_N$.
By (7), (9), and $N\Delta B_N\le2^N$ one obtains $T\le0$.
Thus positive recovery NECESSARILY has
\[
                           N\Delta F_Y(N)>2^N.         \tag{10}
\]
For $K\ge20$, (6) contradicts (10) whenever $N>s_Y$.
Therefore EVERY positive mixed recovery at nonnegative reverse rank has
\[
                              j+1\le s_Y.             \tag{11}
\]
This retains the actual opposite main degree and largest interval.
It does NOT replace them by a first-missing-only linear surrogate.

\subsection{All properly conditioned mixed CLOSED S heads at first missing twenty or above}
At both strict positive windows, the frozen reflection guard gives
$k+1\le j-2$ and $k+1\le a-1$. By (11), the ACTUAL signed opposite
differences satisfy $\Delta q_Y(i)\ge1$ throughout every used head:
the main CI has positive primitive support through $s_Y$, while each
proper shifted-CI summand has nonnegative differences there.
All used coefficients of the raw owner source $R_Z=(2-z)q_Z$ are
nonnegative by the second guard. If $0\le r<k$, its index $r+1$ term
pays $H\ge R_Z(r+1)\Delta q_Y(k-r)\ge C\ge T$.
The frozen relative raw-kernel theorem with $+1$ surplus pays $r\ge k$,
including CLOSED source equality $0\le E_n=E_k$.

For $r\le-2$ positive recovery is impossible. At $r=-1$ the first
source birth pays the unit curvature unless $a\le K-3$, whence
$j\le K-2$ and $k+1\le K-4\le s_Y$; the raw constant term pays.
One-part owners have only this same negative-rank birth exception.
For nonpositive $T$, the usual OLD-A2 head gives $H\ge0$.
This proves the ENTIRE mixed proper CLOSED S comparison for EVERY
$K\ge20$, without capacity, count, weight, padding or amplification bounds.

The remaining mixed first-missing gap is now only $13\le K\le19$
after the accompanying full independent scope/constant verification. The previously frozen
max-twelve criterion, both-forward-rising C resource, both-nonrising
actual theorem and WHOLE closed-S-to-window bridge remain unchanged.
General higher-capacity C still needs proof, and the GENERAL ORIGINAL
conjecture is NOT concluded by this new mixed-source theorem alone.
